\documentclass[11pt,a4paper]{article}
\usepackage[margin=1in]{geometry}
\usepackage{booktabs}
\usepackage{array}
\usepackage{longtable}

\title{IBM BVH Type Mapping}
\author{}
\date{\today}

\begin{document}
\maketitle

This document lists one-to-one type mappings between the original CGAL-based IBM geometry stack and the custom BVH-based implementation, focusing on equivalent functionality.

Current source placement: \texttt{bvh\_defs.h} (data), \texttt{bvh.h} (algorithms), \texttt{eib\_bvh.h} (IBM/BVH bridge and geometry I/O), \texttt{eib\_data.h} (solver SoA), and \texttt{eib\_io.h} (geometry-loading flow).

\renewcommand{\arraystretch}{1.2}
\setlength{\tabcolsep}{6pt}

\begin{longtable}{p{0.19\textwidth} p{0.25\textwidth} p{0.25\textwidth} p{0.25\textwidth}}
\toprule
\textbf{Functionality} & \textbf{CGAL Type} & \textbf{Custom BVH Type} & \textbf{Mapping Notes} \\
\midrule
\endfirsthead

\toprule
\textbf{Functionality} & \textbf{CGAL Type} & \textbf{Custom BVH Type} & \textbf{Mapping Notes} \\
\midrule
\endhead

Point representation & \texttt{Point = Kernel::Point\_2 / Point\_3} & \texttt{Point = GpuArray<Real, AMREX\_SPACEDIM>} & Same semantic role; custom type is GPU-friendly fixed-size storage. \\
\midrule
Vector representation & \texttt{Vector\_CGAL = Kernel::Vector\_2 / Vector\_3} & \texttt{Vec = GpuArray<Real, AMREX\_SPACEDIM>} & Same geometric meaning (direction/difference vectors). \\
\midrule
Inside/outside enum & \texttt{CGAL::Bounded\_side} & \texttt{enum class BoundedSide} & Preserves the same three states: inside, boundary, outside. \\
\midrule
Bounding box type & \texttt{Bbox = CGAL::Bbox\_2 / Bbox\_3} & \texttt{AABB} (alias \texttt{Bbox = AABB}) & Equivalent pruning/containment role for spatial queries. \\
\midrule
2D geometry container & \texttt{GeomType = Polygon2D} (CGAL wrapper) & \texttt{GeomType = Polygon2D} (custom) & Both represent polygon geometry in 2D. \\
\midrule
3D geometry container & \texttt{GeomType = Polyhedron} (\texttt{CGAL::Polyhedron\_3}) & \texttt{GeomType = TriMesh} & Both represent surface meshes; custom stack uses triangle-only explicit indexing. \\
\midrule
Spatial acceleration structure & \texttt{Tree = CGAL::AABB\_tree<Traits>} & \texttt{BVH} & Both accelerate nearest-point and ray-intersection traversal. \\
\midrule
Acceleration internals & \texttt{Primitive + Traits + Tree internals} & \texttt{BVHNode + BVH::nodes} & CGAL hides node internals; custom stack exposes explicit node arrays. \\
\midrule
Closest-point query result & \texttt{Point\_and\_primitive\_id} & \texttt{ClosestPointResult} & Both return closest point and primitive id; custom result also stores distance directly. \\
\midrule
Primitive identifier & \texttt{PrimitiveID = Tree::Primitive\_id} & \texttt{int prim\_id} & Same purpose: identifies edge/face primitive hit by query. \\
\midrule
Element descriptor & \texttt{elm\_descriptor} (iterator or face\_descriptor) & Direct integer indexing & Custom stack replaces handle/descriptor abstractions with linear ids. \\
\midrule
Inside tester functor & \texttt{inside\_t = Side\_of\_triangle\_mesh} (or 2D wrapper) & \texttt{inside\_t = InsideTester} & Unified callable interface retained: \texttt{operator()(Point)}. Current path uses BVH-backed queries in 3D and BVH-accelerated culling for 2D polygon checks. \\
\midrule
Primitive-id map & \texttt{PrimitiveIndexMap = std::map<PrimitiveID, int>} & Not required & Removed because custom pipeline uses direct integer primitive indexing. \\

\bottomrule
\end{longtable}

\end{document}
